由注\ref{rem:relative-open-control-regions}及轨道连续性，
$\{t\in[0,\tau_q):x_q(t)\in\mathcal W\}$ 与
$\{t\in[0,\tau_q):x_q(t)\in\mathcal T\}$
均为 $[0,\tau_q)$ 中的相对开集，因而各自是至多可数个相对开区间的并。
在任一等待区间内，对该区间中的 $t_1<t_2$ 有
\[
x_q(t_2)-x_q(t_1)=\int_{t_1}^{t_2}f_0(x_q(v))\dd v;
\]
在任一完全跟踪区间内，相同等式将 $f_0$ 换为 $f_1$。
这由状态轨道的绝对连续性及上述几乎处处的控制选择得到。
在有限区间端点，通过连续性取极限，积分关系仍成立。
因此，给定每段的起始状态后，由 $f_0,f_1$ 的局部 Lipschitz 性，
该段轨道分别等于对应的唯一自然流或完全跟踪流，直至首次离开相应分区。
这里并未预先假设最优控制只有有限个切换；
各阶段的实际顺序由下面的边界延续与横穿论证确定。
